\begin{abstract}
Post-training changes model behavior, but evaluation alone does not reveal how
those changes are organized in weight space. We study the difference between
initial and final checkpoints, using spectral outliers to locate prominent
directions and matched fine-tuning contrasts to test their behavioral relevance.
Across all $224$ linear maps in the Llama-3-8B Base-to-Instruct delta
\citep{grattafiori2024llama3}, every leading eigenvalue clears a fitted
Marchenko--Pastur visibility edge \citep{marchenko1967}.

The leading attention-output subspace overlaps an independently estimated
refusal direction \citep{arditi2024}. Projecting out its top $128$ directions
changes substring-scored refusal from $98.4\%$ to $14.1\%$, versus $97.7\%$ for
a random subspace, but also causes broad capability loss. Matched harmful and
safe medical fine-tunes \citep{betley2025emergent,turner2025organisms} yield
recurring contrast directions across three 7B/8B families. On
Qwen2.5-Coder-7B \citep{hui2024qwen25coder}, an in-sample projection changes the
measured joint misalignment rate from $2.3\%$ to $0.0\%$, versus $3.4\%$ under
random ablation; a frozen 14B test is inconclusive. In these checkpoints,
spectral geometry exposes behavior-sensitive subspaces, but the capability
damage and failed scale-up criterion do not support an isolated safety mechanism
or prospective causal generalization.
\end{abstract}
